%----------
%	RIESGOS Y MEDIDAS DE MITIGACIÓN
%----------	
La Tabla~\ref{table:risks} recoge los principales riesgos identificados al planificar el trabajo, asociados a los Paquetes de Trabajo (PT) descritos en la Sección~\ref{sec:workmethodology_planning}, junto con su valoración y las medidas de mitigación previstas entonces. Como esta memoria se redacta una vez concluido el trabajo, al análisis inicial le sigue un balance de lo que finalmente ocurrió con cada riesgo y de en qué medida las respuestas previstas resultaron aplicables.


\begin{table}[H]
    \small % Font size
    \ttabbox[\FBwidth]{
        \caption{Riesgos y medidas de mitigación}
        \label{table:risks} % Etiqueta única para la tabla
    }{
        \begin{tabular}{p{4cm} p{2cm} p{2cm} p{1cm} p{4cm}}
            \hline
            \rowcolor{gray!30} % Adding a gray color to the header row
            \textbf{Riesgo} & \textbf{Verosimilitud} & \textbf{Severidad} & \textbf{PTs} & \textbf{Medidas de Mitigación} \\ \hline
            
            \textbf{R1.} Limitaciones de recursos provenientes de IRIS\_IAMEDIA (anotaciones, validación experta, corpus ampliado) & Baja, Media & Media & Todos & Uso del corpus ya validado de IRIS\_IAMEDIA, anotación complementaria propia, generación de ejemplos sintéticos con LLM, retroalimentación puntual en lugar de continua. \\ \hline

            
            \textbf{R2.} Variabilidad y falta de consistencia en los LLM al detectar sesgos & Media & Alta & PT2, PT3, PT4 & Uso de prompting estructurado, definición de casos de test, comparación entre modelos, agente verificador y validación con corpus experto IRIS\_IAMEDIA. \\ \hline

            
            \textbf{R3.} Dificultad para capturar variables cualitativas complejas (androcentrismo, infantilización, etc.) & Media & Media, Alta & PT2, PT4 & Ajuste iterativo de prompts, RAG para aportar contexto, supervisión \textit{human-in-the-loop}, casos anotados por expertas. \\ \hline

            
            \textbf{R4.} Desalineación entre los diagnósticos del sistema y las prácticas y necesidades reales de las redacciones & Media & Media & PT1, PT5 &  Intensificar las consultas con expertas, incorporación de casos de uso reales, test de usabilidad y ajustes iterativos. \\ \hline

            
            \textbf{R5.} Dependencia de herramientas externas o carga académica que afecte al calendario de ejecución & Media & Media & Todos & Planificación semanal, priorización de entregables y alternativas técnicas si fallan las API o servicios externos. \\ \hline

        \end{tabular}
    }
\end{table}

Concluido el trabajo, el balance es el siguiente. Tres de los cinco riesgos se
materializaron, y en ellos las medidas previstas resultaron aplicables, aunque no
siempre suficientes.

El \textbf{R1} se materializó de forma parcial. El corpus anotado estuvo disponible y
bastó para evaluar el plano cuantitativo del sistema, pero no incluía doble
codificación ni marcado experto sobre el texto. Esa carencia impidió calcular la
fiabilidad intercodificadora y validar el plano cualitativo
(Sección~\ref{subsec:limitacion_fiabilidad}). Las medidas previstas de anotación
complementaria y de generación de ejemplos sintéticos no llegaron a emplearse. En su
lugar, la limitación se convirtió en objeto de análisis, ya que el estudio de
sensibilidad por equipo de anotación permitió acotar de forma indirecta el techo de
acuerdo de la tarea (Sección~\ref{subsec:iris_errores}).

El \textbf{R2} se materializó con claridad y acabó siendo uno de los hallazgos del
trabajo, pues los cuatro modelos evaluados mostraron comportamientos marcadamente
distintos entre sí. Se aplicaron las medidas previstas de salida estructurada,
verificación automática de las evidencias, comparación sistemática entre modelos y
validación contra el corpus experto. A ellas se sumó una medida no contemplada
inicialmente, la desactivación del razonamiento explícito en los modelos que lo
admiten, para que la comparación fuese justa
(Sección~\ref{sec:genai_models_iris}).

El \textbf{R3} se materializó y constituye el resultado central del trabajo. La
recuperación aumentada (RAG), prevista como medida de mitigación, se implementó y se
evaluó de forma aislada, con el resultado de que no compensa en este dominio
(Sección~\ref{subsec:iris_descomposicion}). Poder descartarla con evidencia es en sí
mismo una aportación. La supervisión \textit{human-in-the-loop}, en cambio, sí se
incorporó y quedó materializada en la integración del sistema en la herramienta
del proyecto (Sección~\ref{poc}).

El \textbf{R4} no llegó a evaluarse de forma sistemática. La integración del sistema
en la herramienta web del proyecto (Sección~\ref{poc}) y las directrices de uso
editorial que cierran el trabajo (Sección~\ref{sec:conclusions}) responden a la
medida prevista, pero la validación con redacciones reales queda como línea de
trabajo futura (Sección~\ref{sec:futurework_usecases}).

El \textbf{R5} se materializó en dos frentes concretos, y en ambos funcionó la
contingencia prevista. Por un lado, OpenAI anunció la retirada de modelos de su
interfaz, lo que motivó incorporar un modelo más reciente para no depender de uno
solo (Sección~\ref{sec:genai_models_iris}). Por otro, las ejecuciones sobre la granja
de GPU sufrieron interrupciones, para lo que se implementó una ejecución reanudable y
un mecanismo de corte que evita registrar los fallos de infraestructura como
predicciones (Sección~\ref{eval}).


\newpage